\chapter{Field Extensions}

Many equations require a larger field before they can be solved:
\(\mathbb Q(\sqrt2)\) contains a root of \(x^2-2\), and \(\mathbb C\)
contains a root of \(x^2+1\) over \(\mathbb R\).  Extensions also let us
place several related quantities in a single field and, later, enlarge the
scalars on which a vector space is considered.  This chapter makes adjoining
elements precise and establishes the results on which separability and Galois
theory later depend.

\section{Extensions and Degree}\label{sec:extensions-degree}
The vector-space language developed in Chapter~5 is therefore immediately
available: once \(F\subseteq E\), scalar multiplication by elements of \(F\)
makes \(E\) into an \(F\)-vector space, and degree is simply its dimension.
For instance, the common field \(\mathbb Q(\sqrt2,\sqrt3)\) is where the two
square roots can be added and multiplied while still retaining
\(\mathbb Q\) as the base field.

\begin{definition}
An \emph{extension field} \(E/F\), written \(F\subseteq E\), is a field
\(E\) containing \(F\) as a subfield.  Then \(E\) is an \(F\)-vector space.
Its \emph{degree} is
\[
[E:F]=\dim_F E.
\]
The extension is \emph{finite} when this dimension is finite.
\; If
\[
F\subseteq K\subseteq E,
\]
then \(K\) is an \emph{intermediate field} of \(E/F\).
\end{definition}

This definition deliberately reuses familiar linear algebra.  The field
operations in \(E\) are richer than vector-space operations, but the smaller
field \(F\) may still be used as a field of scalars.  Thus the degree measures
how many independent \(F\)-directions must be added in passing from \(F\) to
\(E\).  For example,
\[
[\mathbb Q(\sqrt2):\mathbb Q]=2,
\qquad
[\mathbb C:\mathbb R]=2,
\]
with bases \(1,\sqrt2\) and \(1,i\), respectively.  For elements in a common
ambient field, \(F(S)\) denotes the smallest subfield containing \(F\cup S\).
The inclusion \(\mathbb Q\subseteq\mathbb R\), by contrast, has infinite
degree: no finite list of real numbers spans all of \(\mathbb R\) over
\(\mathbb Q\).  The degree is therefore a genuine measure of the amount of
new scalar information in the extension, not merely a count of symbols.
The infinite extension \(F(t)/F\), formed by adjoining an indeterminate, will
appear in Section~\ref{sec:algebraic-transcendental}; its powers
\(1,t,t^2,\ldots\) cannot fit into any finite
basis.

The Tower Law is the numerical bookkeeping principle for successive
adjunctions.  It says that a basis for the top field is obtained by first
choosing coordinates over the middle field and then expanding each of those
coordinates over the base field.

\begin{theorem}[Tower Law]\label{thm:tower-law}
If
\[
F\subseteq E\subseteq L
\]
and \([E:F]\), \([L:E]\) are finite, then
\[
[L:F]=[L:E][E:F].
\]
\end{theorem}
\begin{proof}
Let \(e_1,\ldots,e_m\) be an \(F\)-basis of \(E\) and
\(u_1,\ldots,u_n\) an \(E\)-basis of \(L\).  The products \(e_i u_j\) span:
expand first in the \(u_j\), then expand each coefficient in the \(e_i\).
If
\[
\sum_{i,j}c_{ij}e_i u_j=0,
\]
then grouping by \(u_j\) and using their independence gives
\[
\sum_i c_{ij}e_i=0
\]
for each \(j\).  Independence of the \(e_i\) makes all \(c_{ij}\) zero.
Thus the products are a basis of size \(mn\).
\end{proof}

\begin{corollary}
If \(F\subseteq E\subseteq L\) and \([L:F]\) is finite, then
\([E:F]\) divides \([L:F]\).  In particular, a prime-degree extension has
no proper intermediate field.
\end{corollary}
\begin{proof}
Apply the Tower Law and factor the prime integer \([L:F]\).
\end{proof}

\section{Algebraic and Transcendental Elements}
\label{sec:algebraic-transcendental}
\begin{definition}
For an extension \(E/F\) and \(\alpha\in E\), the element \(\alpha\) is
\emph{algebraic over \(F\)} if \(f(\alpha)=0\) for some nonzero
\(f\in F[x]\); otherwise it is \emph{transcendental}.  The extension is
\emph{algebraic} when each of its elements is algebraic.
\end{definition}

The evaluation homomorphism
\[
\operatorname{ev}_\alpha:F[x]\longrightarrow E,
\qquad
f\longmapsto f(\alpha)
\]
collects all relations on \(\alpha\).  If \(\alpha\) is algebraic, its
nonzero kernel has a unique monic generator, the \emph{minimal polynomial}
\(m_{\alpha,F}\).

Thus the minimal polynomial is not an arbitrary equation satisfied by
\(\alpha\): it is the irreducible relation of least degree, and it records
exactly the algebraic information lost when the indeterminate \(x\) is
replaced by \(\alpha\).  The next proposition explains why every other
polynomial relation is a consequence of this one.

\begin{proposition}[Minimal-polynomial properties]\label{thm:minimal-polynomial}
The minimal polynomial \(m_{\alpha,F}\) is irreducible and divides every
polynomial in \(F[x]\) which vanishes at \(\alpha\).
\end{proposition}
\begin{proof}
If \(m=gh\), then \(g(\alpha)h(\alpha)=0\), so one factor belongs to the
evaluation kernel.  Minimal positive degree forces the other factor to be a
unit.  For a polynomial \(f\) vanishing at \(\alpha\), divide by \(m\):
\[
f=qm+r,
\qquad
\deg r<\deg m.
\]
Then \(r(\alpha)=0\), so \(r=0\) by minimality.
\end{proof}

\begin{definition}
\(F[\alpha]\) is the subring consisting of polynomial expressions in
\(\alpha\), while \(F(\alpha)\) is the smallest subfield containing \(F\)
and \(\alpha\).
\end{definition}

\begin{theorem}[Simple algebraic extensions]\label{thm:simple-algebraic}
If \(\alpha\) is algebraic and \(d=\deg m_{\alpha,F}\), then
\[
F[\alpha]\cong F[x]/(m_{\alpha,F})=F(\alpha).
\]
The elements
\[
1,\alpha,\ldots,\alpha^{d-1}
\]
are an \(F\)-basis of \(F(\alpha)\), so
\[
[F(\alpha):F]=d.
\]
\end{theorem}
\begin{proof}
Evaluation and the First Isomorphism Theorem give the quotient.  Since the
minimal polynomial is irreducible, that quotient is a field, and therefore
\(F[\alpha]=F(\alpha)\).  Division by the minimal polynomial gives spanning
by the displayed powers.  A relation among them is a smaller nonzero
polynomial in the evaluation kernel, which is impossible.
\end{proof}

\begin{example}[Two square roots]\label{ex:two-square-roots}
The simple-extension theorem now lets us complete the calculation suggested
in Section~\ref{sec:extensions-degree}.  The polynomial \(x^2-3\) has no root in
\(\mathbb Q(\sqrt2)\).  Indeed, every element of this field has the form
\(a+b\sqrt2\) with \(a,b\in\mathbb Q\).  If
\[
(a+b\sqrt2)^2=3,
\]
comparison of the rational and \(\sqrt2\) parts gives
\[
a^2+2b^2=3,
\qquad
2ab=0.
\]
If \(a=0\), then \(b^2=3/2\), impossible in \(\mathbb Q\); if \(b=0\), then
\(a^2=3\), also impossible in \(\mathbb Q\).  Thus \(x^2-3\) is
irreducible over \(\mathbb Q(\sqrt2)\), so it is the minimal polynomial of
\(\sqrt3\) over that field.  The theorem gives
\[
[\mathbb Q(\sqrt2,\sqrt3):\mathbb Q(\sqrt2)]=2.
\]
The Tower Law therefore yields
\[
[\mathbb Q(\sqrt2,\sqrt3):\mathbb Q]
=2\cdot2=4.
\]
This is the first synthesis of irreducibility, minimal polynomials, simple
extension degree, and the Tower Law.  A second square root increases the
degree precisely when it was absent from the first extension.
\end{example}

\begin{proposition}[The transcendental case]
If \(t\) is transcendental over \(F\), then
\[
F[t]\cong F[x],
\qquad
F(t)\cong F(x).
\]
In particular \(F(t)/F\) is infinite.
\end{proposition}
\begin{proof}
The evaluation kernel is zero.  Passing from the polynomial ring to its
field of fractions gives the second isomorphism, and the powers of \(t\) are
linearly independent over \(F\).
\end{proof}

The algebraic and transcendental cases have opposite shapes.  An algebraic
element has a power that can be reduced to a combination of finitely many
earlier powers, whereas the powers of a transcendental element never satisfy
such a relation.  Passing from \(F[t]\) to \(F(t)\) adds fractions, just as
passing from \(\mathbb Z\) to \(\mathbb Q\) adds quotients; it does not create
a polynomial relation for \(t\).
\[
\begin{array}{c|c|c}
& \text{algebraic }\alpha & \text{transcendental }t\\ \hline
\text{polynomial relations} & \text{generated by }m_{\alpha,F} & \text{none}\\
\text{associated field} & F(\alpha)\cong F[x]/(m_{\alpha,F}) & F(t)\cong F(x)\\
\text{degree over }F & \text{finite} & \text{infinite}
\end{array}
\]

\begin{example}
Eisenstein's criterion makes \(x^3-2\) the minimal polynomial of
\(\sqrt[3]2\) over \(\mathbb Q\).  Thus
\[
\mathbb Q(\sqrt[3]2)
=\mathbb Q\oplus\mathbb Q\sqrt[3]2\oplus\mathbb Q\sqrt[3]4.
\]
In contrast, \(F(x)\) contains rational functions not in \(F[x]\).
\end{example}

\begin{definition}
An extension is \emph{finitely generated} if it has the form
\[
F(\alpha_1,\ldots,\alpha_r).
\]
\end{definition}

\begin{proposition}\label{prop:algebraic-towers}
Finite extensions are algebraic, finitely generated algebraic extensions are
finite, and algebraicity is transitive through towers.
\end{proposition}
\begin{proof}
First suppose \([E:F]=n\), and let \(a\in E\).  The \(n+1\) elements
\[
1,a,\ldots,a^n
\]
are linearly dependent over \(F\).  Their dependence relation is a nonzero
polynomial over \(F\) which kills \(a\), proving that \(E/F\) is algebraic.

Next let
\[
E=F(\alpha_1,\ldots,\alpha_r)
\]
with every \(\alpha_i\) algebraic over \(F\).  Each \(\alpha_i\) remains
algebraic over the intermediate field already constructed, because its
polynomial over \(F\) is also a polynomial over that field.  Thus every
successive extension is finite by Theorem~\ref{thm:simple-algebraic}; repeated
use of the Tower Law makes \(E/F\) finite.

Finally suppose \(F\subseteq E\subseteq L\), \(E/F\) is algebraic, and
\(\alpha\in L\) is algebraic over \(E\).  Write a monic relation
\[
\alpha^n+a_{n-1}\alpha^{n-1}+\cdots+a_0=0,
\qquad
a_i\in E,
\]
and put
\[
K=F(a_0,\ldots,a_{n-1}).
\]
The elements \(a_i\) are algebraic over \(F\), so \(K/F\) is finite by the
second part.  The displayed relation shows that \(\alpha\) is algebraic over
\(K\); hence \(K(\alpha)/K\) is finite.  The Tower Law gives
\[
[K(\alpha):F]<\infty.
\]
Its first part now says that \(\alpha\) is algebraic over \(F\).  This proves
transitivity.
\end{proof}

\section{Root Adjunction and Splitting Fields}
When an irreducible polynomial has no root in the current field, this is not
an obstruction to doing algebra with a root.  The quotient construction below
creates the smallest possible field in which a formal symbol is required to
satisfy that polynomial.

\begin{theorem}[Root adjunction]\label{thm:root-adjunction}
If \(f\in F[x]\) is irreducible, then
\[
F[x]/(f)
\]
is a field containing \(F\), and \(x+(f)\) is a root of \(f\).  Its minimal
polynomial is \(f\), and its degree over \(F\) is \(\deg f\).
\end{theorem}
\begin{proof}
The quotient is the formal way to impose the single relation \(f(x)=0\).
Because \(f\) is irreducible, the
\hyperref[thm:irreducible-maximal-field]{irreducible-quotient theorem of
Chapter~4} makes the principal ideal \((f)\) maximal; hence the quotient is a
field rather than merely a ring.
Constants embed as their residue classes, and the class
\[
\alpha=x+(f)
\]
is the formal adjoined root since \(f(\alpha)=f(x)+(f)=0\).  The evaluation
kernel is \((f)\); the remaining claims follow from
Theorem~\ref{thm:simple-algebraic}.
\end{proof}

\begin{example}[A root made concrete]
Taking \(f=x^2-2\in\mathbb Q[x]\) gives
\[
\mathbb Q[x]/(x^2-2).
\]
If \(\alpha=x+(x^2-2)\), then \(\alpha^2=2\), and every residue class has
a unique form
\[
a+b\alpha,
\qquad a,b\in\mathbb Q.
\]
For instance,
\[
(a+b\alpha)(c+d\alpha)=(ac+2bd)+(ad+bc)\alpha.
\]
Sending \(\alpha\) to the usual real number \(\sqrt2\) identifies this
quotient with \(\mathbb Q(\sqrt2)\).  The quotient construction is useful
because it works before a convenient ambient field, such as \(\mathbb R\),
has been chosen.
\end{example}

\begin{lemma}[Extension of embeddings]\label{lem:root-extension}
Let \(\sigma:K\to\Omega\) be an \(F\)-embedding.  Let \(\alpha\) be
algebraic over \(K\), with minimal polynomial \(m\).  Extensions
\[
\widetilde\sigma:K(\alpha)\longrightarrow\Omega
\]
of \(\sigma\) correspond bijectively to roots in \(\Omega\) of the
coefficientwise image \(\sigma(m)\).
\end{lemma}
\begin{proof}
Write
\[
m(x)=a_0+a_1x+\cdots+a_nx^n,
\qquad
\sigma(m)(x)=\sigma(a_0)+\sigma(a_1)x+\cdots+\sigma(a_n)x^n.
\]
For a root \(\beta\) of this transported polynomial, the map
\[
\sum a_ix^i\longmapsto\sum\sigma(a_i)\beta^i
\]
has \(m\) in its kernel and descends to a homomorphism from
\[
K[x]/(m)\cong K(\alpha)
\]
to \(\Omega\).  This is well-defined precisely because the defining relation
\(m(\alpha)=0\) is respected by the chosen root \(\beta\).  The source is a
field, so it is injective.  It is uniquely determined by its values on \(K\)
and \(\alpha\).  Conversely an extension
sends \(m(\alpha)=0\) to
\[
\sigma(m)(\widetilde\sigma(\alpha))=0.
\]
\end{proof}

\begin{definition}
A \emph{splitting field} of a nonconstant polynomial \(f\in F[x]\) is an
extension generated over \(F\) by the roots of \(f\), in which \(f\) splits
into linear factors.
\end{definition}

\begin{theorem}[Splitting fields]\label{thm:splitting-existence}
Every nonconstant polynomial in \(F[x]\) has a splitting field, unique up to
an \(F\)-isomorphism.
\end{theorem}
\begin{proof}
For existence, induct on \(n=\deg f\).  There is nothing to prove for
\(n=1\).  Otherwise choose an irreducible factor \(p\) of \(f\), and use
Theorem~\ref{thm:root-adjunction} to obtain a field \(K/F\) containing a root
\(\alpha\) of \(p\).  In \(K[x]\), division gives
\[
f(x)=(x-\alpha)g(x),
\qquad
\deg g=n-1.
\]
By induction, \(g\) has a splitting field \(L/K\).  Then \(f\) splits in
\(L\), and \(L\) is generated over \(F\) by \(\alpha\) together with the
roots of \(g\), which are precisely all roots of \(f\).

For uniqueness, let \(L=F(\alpha_1,\ldots,\alpha_r)\) and \(L'\) be
splitting fields of \(f\).  Start with \(1_F:F\to L'\).  Suppose that an
embedding has been defined on \(F(\alpha_1,\ldots,\alpha_{j-1})\).  The
minimal polynomial of \(\alpha_j\) over this field divides \(f\), after the
coefficients have been transported.  Since \(f\) splits in \(L'\), that
transported minimal polynomial has a root in \(L'\).  Lemma
\ref{lem:root-extension} extends the embedding over \(\alpha_j\).

After the last step we have an \(F\)-embedding \(\sigma:L\to L'\).  If
\[
f(x)=c\prod_i(x-\alpha_i)
\]
in \(L[x]\), then applying \(\sigma\), which fixes the coefficients of
\(f\), gives
\[
f(x)=c\prod_i(x-\sigma(\alpha_i))
\]
in \(L'[x]\).  Thus the images are the complete roots of \(f\), with
multiplicities.  Since \(L'\) is generated by those roots, it lies in
\(\sigma(L)\).  Therefore \(\sigma\) is onto and is an isomorphism.
\end{proof}

\begin{example}
The splitting field of \(x^2-2\) is \(\mathbb Q(\sqrt2)\).  The three roots
of \(x^3-2\) are
\[
\sqrt[3]2,\qquad\zeta_3\sqrt[3]2,\qquad\zeta_3^2\sqrt[3]2,
\]
so its splitting field is
\[
\mathbb Q(\sqrt[3]2,\zeta_3).
\]
The first adjunction supplies one real root, but it does not supply the two
conjugate nonreal roots.  The primitive cube root \(\zeta_3\) is precisely
what rotates \(\sqrt[3]2\) into those roots, so adjoining it produces the
whole splitting field.  This is the first indication that the symmetries of
all roots together, rather than a single chosen root, are what Galois theory
will study.  The polynomial \((x-1)(x+2)\) already splits over
\(\mathbb Q\), so its splitting field is \(\mathbb Q\).  In characteristic
\(p\), the polynomial \(x^p-x\) splits in \(\mathbb F_p\), because each
\(a\in\mathbb F_p\) satisfies \(a^p=a\), and it has degree \(p\).
\end{example}

\section{Algebraically Closed Fields and Embeddings}
\begin{definition}
A field \(\Omega\) is \emph{algebraically closed} if every nonconstant
polynomial in \(\Omega[x]\) has a root in \(\Omega\).
\end{definition}
\begin{proposition}
A field is algebraically closed if and only if every nonconstant polynomial
splits into linear factors.
\end{proposition}
\begin{proof}
Repeatedly factor out a root and induct on degree.  The converse is immediate.
\end{proof}

\begin{definition}
An \emph{\(F\)-embedding} \(\sigma:E\to\Omega\) is a field homomorphism
between extensions of \(F\) that fixes \(F\) pointwise.  Its images of an
algebraic element are called its \emph{conjugates}.
\end{definition}

\begin{proposition}[Embeddings and conjugates]\label{prop:embeddings-roots}
If \(\Omega\) is algebraically closed and \(\alpha\) is algebraic over \(F\),
then
\[
\sigma\longmapsto\sigma(\alpha)
\]
is a bijection from \(F\)-embeddings
\[
F(\alpha)\longrightarrow\Omega
\]
to roots of \(m_{\alpha,F}\) in \(\Omega\).
\end{proposition}
\begin{proof}
An embedding sends the equation \(m_{\alpha,F}(\alpha)=0\) to the required
root relation.  Thus every image is a root of the minimal polynomial; these
are exactly the possible algebraic realizations of \(\alpha\) compatible with
the base field.  Conversely, a root gives the unique embedding by
Lemma~\ref{lem:root-extension}.  The two constructions are inverse because
an embedding of \(F(\alpha)\) is determined by its value on \(\alpha\).
\end{proof}

\begin{example}
The conjugates of \(\sqrt2\) over \(\mathbb Q\) are
\[
\sqrt2,\qquad-\sqrt2.
\]
There are correspondingly two \(\mathbb Q\)-embeddings of
\(\mathbb Q(\sqrt2)\) into \(\mathbb C\).  The conjugates of
\(\sqrt[3]2\) are its three complex cube roots; Chapter~10 will explain why
the number of distinct conjugates can be smaller than the degree in positive
characteristic.
\end{example}

More explicitly, the three \(\mathbb Q\)-embeddings of
\(\mathbb Q(\sqrt[3]2)\) into \(\mathbb C\) are determined by the three
possible images
\[
\sqrt[3]2,
\qquad
\zeta_3\sqrt[3]2,
\qquad
\zeta_3^2\sqrt[3]2.
\]
Only the first has image equal to the original real cubic field.  The other
two land in different conjugate subfields of its splitting field.  This is
why embeddings are the right preliminary notion: automorphisms are the
special embeddings whose image remains in the same field.

The next extension theorem is the standard place where Zorn's Lemma enters
field theory.  Start with every partial extension of the prescribed embedding,
order these choices by further extension, take the union along a chain, and
observe that a maximal partial embedding can be enlarged unless its domain is
already all of \(E\).  The proof records the details because this pattern will
return when algebraic closures and normal extensions are compared.

\begin{theorem}[Extension of embeddings]\label{thm:embedding-extension}
If \(E/K\) is algebraic, \(\Omega\) is algebraically closed, and
\(\sigma:K\to\Omega\) is an \(F\)-embedding, then \(\sigma\) extends to an
\(F\)-embedding \(E\to\Omega\).
\end{theorem}
\begin{proof}
Let \(\mathcal P\) consist of pairs \((M,\tau)\) with
\[
K\subseteq M\subseteq E
\]
and \(\tau:M\to\Omega\) an \(F\)-embedding extending \(\sigma\).  Define
\[
(M,\tau)\leq(N,\rho)
\]
when \(M\subseteq N\) and \(\rho|_M=\tau\).  This is a set: the possible
domains are subsets of the fixed set \(E\), and possible maps are subsets of
\(E\times\Omega\).

Let \(\mathcal C\) be a chain.  Its domains have union
\[
M_{\mathcal C}=\bigcup_{(M,\tau)\in\mathcal C}M.
\]
For two elements of this union, one chain member contains both, so it
contains their sum, product, and inverses.  Hence the union is a field.  The
maps agree on overlaps because the pairs are comparable; their union is
therefore a well-defined \(F\)-embedding
\(\tau_{\mathcal C}:M_{\mathcal C}\to\Omega\).  It is an upper bound.
Zorn's Lemma supplies a maximal pair \((M,\tau)\).

If \(M\ne E\), choose \(\alpha\in E\setminus M\).  Since \(E/M\) is
algebraic, \(\alpha\) has a minimal polynomial over \(M\).  The coefficientwise
image of that polynomial has a root in the algebraically closed field
\(\Omega\), and Lemma~\ref{lem:root-extension} extends \(\tau\) to
\(M(\alpha)\), contradicting maximality.  Thus \(M=E\).
\end{proof}

\section{Algebraic Closures}
\begin{definition}
An \emph{algebraic closure} \(\overline F/F\) is an algebraic extension
\(\overline F\) that is algebraically closed.
\end{definition}

An algebraic closure is a sufficiently large, but still algebraic, universe
in which every polynomial over the starting field can be completely split.
The construction below first enlarges a field so that every current polynomial
has a root, then repeats, and finally retains the elements algebraic over the
original field.  That final restriction is what keeps the closure algebraic
over \(F\).

\begin{theorem}[Existence of algebraic closures]\label{thm:algebraic-closure-existence}
Every field has an algebraic closure.
\end{theorem}
\begin{proof}
We use a construction that works within ordinary sets throughout.  For any
field \(K\), let \(S_K\) be the set of nonconstant polynomials in \(K[x]\).
Introduce one indeterminate \(X_f\) for each \(f\in S_K\), and form
\[
R_K=K[X_f:f\in S_K].
\]
Let \(I_K\) be the ideal generated by the elements
\[
f(X_f),
\qquad
f\in S_K.
\]
This ideal is proper.  Indeed, any alleged relation \(1\in I_K\) involves
only finitely many variables \(X_{f_1},\ldots,X_{f_r}\).  Adjoin roots one at
a time by Theorem~\ref{thm:root-adjunction}; in the resulting field every
\(f_i(X_{f_i})\) vanishes, while \(1\) does not.  Thus no such relation can
exist.

By the \hyperref[thm:maximal-ideal-existence]{maximal-ideal theorem}, choose
a maximal ideal \(\mathfrak m_K\supseteq I_K\), and set
\[
K^+=R_K/\mathfrak m_K.
\]
The quotient is a field containing a copy of \(K\): the ideal
\(\mathfrak m_K\cap K\) is an ideal of the field \(K\), and it cannot equal
\(K\), since then \(1\in\mathfrak m_K\).  Hence
\[
\mathfrak m_K\cap K=(0),
\]
so the composite \(K\to R_K\to K^+\) is injective.  Every nonconstant
polynomial \(f\in K[x]\) has the root \(X_f+\mathfrak m_K\) in \(K^+\).

Starting with \(F_0=F\), recursively construct
\[
F_0\subseteq F_1\subseteq F_2\subseteq\cdots,
\qquad
F_{n+1}=F_n^+,
\]
and set
\[
\Omega=\bigcup_{n\geq0}F_n.
\]
This union is a field.  If \(g\in\Omega[x]\) is nonconstant, its finitely
many coefficients lie in some \(F_n\), and then \(g\) has a root in
\(F_{n+1}\subseteq\Omega\).  Hence \(\Omega\) is algebraically closed.

Finally let \(\overline F\) be the subfield of \(\Omega\) consisting of
elements algebraic over \(F\).  This set is a field: if \(\alpha,\beta\) are
algebraic over \(F\), then \(F(\alpha,\beta)/F\) is finite.  Hence
\[
\alpha+\beta,\qquad \alpha-\beta,\qquad \alpha\beta,
\qquad \alpha^{-1}\quad(\alpha\ne0)
\]
belong to a finite extension of \(F\), and are therefore algebraic over \(F\).
If \(h\in\overline F[x]\) is nonconstant, its coefficients lie in a finite
algebraic extension \(K/F\) inside \(\overline F\).  A root \(b\in\Omega\)
of \(h\) is algebraic over \(K\), hence algebraic over \(F\); so
\(b\in\overline F\).  Therefore \(\overline F\) is algebraically closed and
is an algebraic closure of \(F\).
\end{proof}

\begin{theorem}[Uniqueness of algebraic closures]
Any two algebraic closures of \(F\) are \(F\)-isomorphic.
\end{theorem}
\begin{proof}
For closures \(\overline F,\overline F'\), Theorem
\ref{thm:embedding-extension} extends \(1_F\) to an embedding
\[
\sigma:\overline F\longrightarrow\overline F'.
\]
The image is algebraically closed: if a polynomial has coefficients in
\(\sigma(\overline F)\), transport its coefficients back through \(\sigma\),
take a root in \(\overline F\), and transport that root forward.  Now take
\(\beta\in\overline F'\).  It is algebraic over \(\sigma(\overline F)\).
Its irreducible minimal polynomial has a root in the algebraically closed
field \(\sigma(\overline F)\), so it is linear.  Hence
\(\beta\in\sigma(\overline F)\).  Therefore the image is all of
\(\overline F'\).
\end{proof}

\begin{remark}
An algebraic closure is thus unique only up to an isomorphism fixing \(F\).
There is usually no distinguished isomorphism between two constructions, so one fixes
a chosen algebraic closure when it is useful to speak of all conjugates in a
single ambient field.
\end{remark}

\section*{Exercises}
\begin{exercise}[Basic] Prove directly that
\[
[\mathbb Q(\sqrt2,\sqrt3):\mathbb Q]=4.
\]
\end{exercise}
\begin{exercise}[Basic] Find the minimal polynomial of \(1+\sqrt2\) over
\(\mathbb Q\).
\end{exercise}
\begin{exercise}[Basic] Show that \(\mathbb Q(\sqrt2)=\mathbb Q(1+\sqrt2)\).
\end{exercise}
\begin{exercise}[Basic] Prove that \(\mathbb Q(\sqrt2)\cap\mathbb Q(\sqrt3)
=\mathbb Q\).
\end{exercise}
\begin{exercise}[Basic] Show that \(F[t]\) is not a field when \(t\) is
transcendental over \(F\), whereas \(F(t)\) is a field.
\end{exercise}
\begin{exercise}[Core] Show that if \(\alpha\) has odd degree over \(F\),
then \(F(\alpha^2)=F(\alpha)\).
\end{exercise}
\begin{exercise}[Core] Let \(\alpha\) have minimal polynomial
\(x^4+2x^2+2\) over \(F\).  Write \(\alpha^7\) as an \(F\)-linear combination
of \(1,\alpha,\alpha^2,\alpha^3\).
\end{exercise}
\begin{exercise}[Core] Prove that every element of a finite extension
\(E/F\) has minimal polynomial of degree at most \([E:F]\).
\end{exercise}
\begin{exercise}[Core] If \(F\subseteq E\subseteq L\), with
\([L:F]=12\), list the possible degrees of \(E/F\).
\end{exercise}
\begin{exercise}[Core] Construct the splitting field of \(x^4-2\) over
\(\mathbb Q\) and determine its degree.
\end{exercise}
\begin{exercise}[Core] Prove that an intermediate field of a degree-six
extension cannot have degree four over the base.
\end{exercise}
\begin{exercise}[Further] Show that the number of \(F\)-embeddings of
\(F(\alpha)\) into an algebraic closure is at most \(\deg m_{\alpha,F}\).
\end{exercise}
\begin{exercise}[Further] Let \(f\in F[x]\) be irreducible and let
\(\alpha,\beta\) be roots of \(f\) in an algebraically closed field.  Use the
root-extension lemma to construct an \(F\)-isomorphism
\[
F(\alpha)\longrightarrow F(\beta).
\]
\end{exercise}
\begin{exercise}[Further] Prove that a field which is algebraic over an
algebraically closed field is that field itself.
\end{exercise}
\begin{exercise}[Further] Show that \(\mathbb C\) is an algebraic closure of
\(\mathbb R\), using the Fundamental Theorem of Algebra as an external input.
A Galois-theoretic proof of that theorem appears later in Chapter~11.
\end{exercise}
\begin{exercise}[Looking Ahead] Show that every finite extension of a finite
field is a splitting field of a polynomial over that finite field.
\end{exercise}
\begin{exercise}[Looking Ahead] Let \(f\in F[x]\) be irreducible.  Show that
every \(F\)-embedding of \(F(\alpha)\) into an algebraic closure sends
\(\alpha\) to a root of \(f\), and compare this with the automorphisms of a
splitting field of \(f\).
\end{exercise}
